\documentclass{article}

\PassOptionsToPackage{numbers,sort&compress}{natbib}
\usepackage[dblblindworkshop]{neurips_2026}
\workshoptitle{New in Machine Learning (NewInML) at NeurIPS 2026}

\usepackage[utf8]{inputenc}
\usepackage[T1]{fontenc}
\usepackage{hyperref}
\usepackage{url}
\usepackage{booktabs}
\usepackage{amsmath}
\usepackage{graphicx}
\usepackage{xcolor}
\usepackage{microtype}

\title{HydraDG: Governed Context Interventions with Failure-Complete Custody for Agent Experiments}

\author{Anonymous Author(s)}

\begin{document}

\maketitle

\begin{abstract}
Agent systems increasingly rely on context interventions---structured retrieval, memory graphs, and prompt contracts---yet published evaluations often omit null outcomes, parser failures, and negative results.
We present HydraDG, a governed experimental framework that binds every substantive model invocation, tool call, and handoff to a Failure-Complete Object (FCO) and Failure-Complete Graph (FCG) custody layer.
HydraDG preregisters hypotheses, freezes case manifests and scorers, records raw prompt and response hashes, and preserves terminal states including timeouts, abstentions, and malformed outputs without silent substitution.
We report two preregistered confirmatory experiments (EXP-008, EXP-009) comparing flat prose versus structured FCG retrieval for failure-prevention endpoints; both terminate as \emph{underpowered} under bounded case replication, establishing that confirmatory effect claims are not supported while retaining full negative and null evidence in custody.
We argue that custody-first evaluation infrastructure is a prerequisite for reproducible agent science rather than an optional logging detail.
\end{abstract}

\section{Introduction}

Large language model agents are evaluated on task success rates, benchmark scores, and human preference.
Less often preserved are the conditions under which experiments fail: swapped models, stale runtime selection, partial matrix execution, or JSON outputs that never parse.
When negative and null cells disappear from the public record, later work cannot distinguish \emph{no effect} from \emph{no measurement}.

HydraDG addresses this gap with a custody-first experimental contract.
Every actor---human operator, frontier agent, local model runtime, or deterministic script---must emit or be represented by a handoff receipt before its output is promoted.
Receipts bind actor identity, host, repository state, input hashes, raw output hashes, evidence class, claim ceiling, and signature state.
Probabilistic model outputs are never treated as authority; they are evidence objects subject to deterministic scoring and graph append rules.

This paper describes the framework architecture and reports terminal results from the IC Failure Learning program's structured-retrieval falsification lane.
We do \emph{not} claim that structured FCG context improves failure prevention; preregistered experiments EXP-008 and EXP-009 closed as underpowered.
We also do not treat exploratory secondary patterns as promoted findings.

\section{Related Work}

\textbf{Agent evaluation.} Benchmark suites measure capability but rarely require publication of failed runs, malformed generations, or infrastructure blocks.
HydraDG complements benchmarks with mandatory negative capture.

\textbf{Provenance and reproducibility.} Scientific workflows use preregistration, version pinning, and artifact hashing; agent stacks often stop at prompt logging.
FCO/FCG extends content-addressed lineage to cross-actor handoffs so that a stale agent cannot write after ownership moves.

\textbf{Structured context and retrieval.} Graph-augmented prompting and memory systems assume structured context helps; falsification requires frozen contracts and power-aware verdicts.
HydraDG enforces both.

\section{Framework}

\subsection{Custody objects}

An \textbf{FCO} materializes one bounded transformation: e.g., a model call, scorer pass, or ingest operation.
An \textbf{FCG} append records directed edges between FCOs with hash-linked roots.
Claim ceilings (exploratory mechanistic falsification, descriptive only, etc.) are declared at preregistration and enforced at promotion time.

\subsection{Experimental freeze}

Each experiment declares:
\begin{itemize}
  \item frozen case manifest and prompt-contract hashes;
  \item conditions (e.g., flat prose vs.\ structured FCG);
  \item models with expected runtime digests;
  \item primary endpoint and aggregation rule;
  \item explicit gates for confirmatory vs.\ exploratory claims.
\end{itemize}
Execution emits \texttt{START\_RECEIPT}, per-cell raw outputs with prompt/response SHA-256, parser state, and terminal receipts.
Integrity failure (duplicate keys, missing dependencies, silent model substitution) blocks promotion; scientific failure (timeout, abstention, malformed JSON) is retained as terminal evidence.

\subsection{Governed local execution}

Where preregistered, local model calls route through a governed bridge that records selection age, swap posture, and resolved digest.
Direct runtime invocation and governed invocation are separate evidence classes; one is not silently labeled as the other.

\section{Experimental Setup}

\subsection{Task family}

The IC Failure Learning suite uses synthetic agent-failure vignettes with preregistered endpoints such as E06 (prevents exposure of out-of-vault media).
Cases are drawn from a frozen \texttt{CASES.jsonl} manifest; each case receives multiple replicates per condition.

\subsection{Conditions and models}

\textbf{C0 (FLAT\_PROSE):} failure-handling guidance as unstructured text.\\
\textbf{C1 (STRUCTURED\_FCG):} the same semantic content composed through a frozen FCG retrieval contract.

EXP-008 and EXP-009 used local models \texttt{qwen3:1.7b} and \texttt{qwen2.5-coder:7b} under identical case and scorer contracts.
Thinking/reasoning modes were held consistent with preregistration.
Primary aggregation: majority of three replicates per case.

\subsection{AI and agent methodology disclosure}

Frontier coding agents assisted with repository tooling and manuscript preparation; they are not listed as authors.
All scientific endpoints come from preregistered local model executions with hashed prompts and responses.
Routine editing does not constitute a reported experimental intervention.

\subsection{Power and claim gates}

E06 endpoints are known limited-power under the frozen case set.
Confirmatory ordering claims require preregistered discordant pairs; exploratory secondary statistics are quarantined from primary promotion.

\section{Results}

Table~\ref{tab:terminal} summarizes terminal verdicts for preregistered confirmatory experiments.
Both experiments executed the full raw matrix ($n=300$ cells each) with complete custody append.

\begin{table}[t]
  \centering
  \caption{Terminal confirmatory experiments (HydraDG IC Failure Learning). \emph{Underpowered} denotes insufficient paired evidence for confirmatory promotion, not proof of null effect.}
  \label{tab:terminal}
  \begin{tabular}{llll}
    \toprule
    Experiment & Primary verdict & Raw cells & Valid parse rate \\
    \midrule
    EXP-008 & UNDERPOWERED & 300 & 0.907 \\
    EXP-009 & UNDERPOWERED & 300 & 0.883 \\
    \bottomrule
  \end{tabular}
\end{table}

\textbf{EXP-008.} Structured FCG retrieval did not yield promotable confirmatory evidence on E06 under the frozen design ($n=2$ paired cases per model stratum at aggregation scope). Data quality metrics show non-zero malformed and abstention rates preserved in custody.

\textbf{EXP-009.} Primary verdict remains UNDERPOWERED for confirmatory E06 ordering. An exploratory secondary pattern was recorded as directionally positive at the descriptive layer only; it is \textbf{not promoted} to a causal claim and does not motivate redesign of the primary hypothesis.

\textbf{Stage-2 baseline.} Prior IC Failure Learning stage-2 execution (414 scored rows across two models) established \texttt{FAILURE\_LEARNING\_BEHAVIOR\_IMPROVEMENT\_NOT\_ESTABLISHED} with committed post-model MMR state, providing historical context for the structured-retrieval falsification lane.

\section{Discussion}

HydraDG demonstrates that agent experiments can be run with the same negative-result integrity expected in preclinical trials: failures are first-class, hashes are mandatory, and claim ceilings block over-interpretation.
The EXP-008/009 underpowered terminations are scientifically informative---they bound what cannot yet be claimed---while preserving malformed and timeout cells for audit.

Limitations include bounded case replication, local-only model lanes in the reported confirmatory matrix, and absence of runtime-equivalent successor replication in this submission's primary results.
Future work includes completing preregistered successor lanes under resource gates without merging partial counts into confirmatory tables.

\section{Conclusion}

We presented HydraDG, a custody-first framework for agent context-intervention experiments, and reported two terminal underpowered confirmatory studies that falsify premature promotion of structured FCG retrieval effects on failure-prevention endpoints.
We recommend adoption of FCO/FCG custody, explicit claim ceilings, and terminal-state preservation as baseline infrastructure for NewInML-style agent evaluations.

\begin{thebibliography}{9}

\bibitem{prereg2026}
Anonymous. (2026). HydraDG IC Failure Learning preregistrations EXP-008 and EXP-009. Internal frozen artifacts; anonymized submission.

\bibitem{stage2}
Anonymous. (2026). IC Failure Learning Stage-2 closeout receipt. Internal frozen artifacts; anonymized submission.

\bibitem{neurips2026}
NeurIPS. (2026). New in Machine Learning (NewInML) workshop call for papers. Non-archival workshop track.

\end{thebibliography}

\end{document}
